\section{Preliminaries}
\label{sec:preliminaries}

This section introduces the fundamental concepts and definitions required to
understand the remainder of this report. It establishes the formal game-theoretic
model of chess, derives the associated search problem, and motivates the use of
heuristic evaluation and time-constrained search. All concepts presented in this
section are standard in the field of computer chess and provide a common basis
for the subsequent chapters.

\subsection{Chess as a Game-Theoretic Model}

Chess is modeled as a deterministic, two-player, zero-sum game with perfect
information. At any point in time, the complete game state is fully observable
to both players, and no randomness is involved in the game mechanics.

A game state consists of a board configuration together with additional state
information such as the side to move, castling rights, and en-passant
availability. Players alternate turns by selecting legal moves, and the game
terminates in a win, loss, or draw according to the rules of chess.

\paragraph{Formal Representation}
From a theoretical perspective, the game can be represented as a directed graph.
Each node corresponds to a legal position, and each directed edge corresponds to
a legal move transforming one position into another. Terminal nodes represent
positions in which no legal moves are available, such as checkmate or stalemate,
as well as positions that are declared drawn by the rules of chess.

\paragraph{Game Tree Interpretation}
For decision making, the directed graph induces a game tree rooted at the
current position. Each level of the tree corresponds to one ply. Due to the high
branching factor of chess, which typically lies between 30 and 40 in middlegame
positions, the size of the game tree grows exponentially with depth. As a
result, even moderately deep exhaustive search is computationally infeasible.

\subsection{Search Problem Formulation}

Given a current position, the task of a chess engine is to select a legal move
that maximizes the expected outcome under the assumption of optimal play by both
sides. In the classical formulation, outcomes are mapped to numerical values,
for example win, draw, and loss.

This decision problem is commonly formulated as a game tree search problem, in
which the engine explores a subset of the game tree to estimate the value of
available moves. The root of the tree corresponds to the current position, while
internal nodes represent positions reachable through sequences of legal moves.

\paragraph{Resource Constraints}
Due to the exponential growth of the search tree, it is impossible to evaluate
all reachable positions. Consequently, the search must be limited by depth,
time, or other computational resources. This limitation introduces approximation:
the engine can no longer guarantee optimal play, but instead aims to select the
best move according to the information explored within the available budget.

\paragraph{Role of Approximation}
The quality of a move selected by a chess engine therefore depends on two main
factors: the effectiveness of the search algorithm in selecting which parts of
the tree to explore, and the accuracy of the evaluation function used to estimate
the value of positions at the search frontier.

\subsection{Minimax, Negamax, Quiescene and Alpha-Beta Pruning}
\label{subsec:minimax_alphabeta}

The classical algorithmic solution to game tree search in two-player zero-sum
games with perfect information is the Minimax algorithm. Minimax assumes that
one player aims to maximize the evaluation value of a position, while the
opponent aims to minimize it. Terminal positions are assigned fixed values
corresponding to win, loss, or draw, while non-terminal leaf positions are
evaluated heuristically.

The algorithm explores the game tree recursively and propagates values from the
leaves toward the root by alternating between maximization and minimization
steps. A detailed introduction to Minimax can be found in standard AI textbooks,
such as Russell and Norvig~\cite{russell2010aima}.

\paragraph{Negamax Formulation}
In zero-sum games, the Minimax algorithm can be simplified using the Negamax
formulation. Instead of explicitly distinguishing between maximizing and
minimizing players, Negamax expresses all evaluations from the perspective of the
side to move. The value of a position is defined as the negation of the value
obtained by the opponent after making a move.

This formulation exploits the inherent symmetry of zero-sum games and allows the
search to be implemented using a single recursive procedure. Negamax is
functionally equivalent to Minimax but leads to more compact and uniform code,
which is why it is widely used in practical chess engines.

\paragraph{Quiescene Search}
Depth-limited game tree search combined with static evaluation functions suffers
from a well-known problem referred to as the \emph{horizon effect}. Tactical
events such as captures, checks, or promotions may occur just beyond the search
depth and are therefore not accounted for by the evaluation function. As a
result, the engine may incorrectly assess positions as stable even though
significant tactical changes are imminent.

Quiescence search is a standard technique to solve this problem. Instead of
evaluating a position immediately when the nominal search depth is reached, the
search is selectively extended in positions that are considered tactically
unstable. Typically, only forcing moves such as captures or checking moves are
explored, while quiet moves are excluded.

The goal of quiescence search is to reach a \emph{quiescent} position, that is, a
position in which no immediate tactical exchanges are available and the static
evaluation is therefore more reliable. Once such a position is reached, the
evaluation function is applied and the resulting score is propagated back up the
search tree.

In practice, quiescence search is commonly implemented using Alpha-Beta pruning
with a so-called \emph{stand-pat} evaluation, which represents the static
evaluation of the current position before any tactical extensions are explored.
If the stand-pat score already exceeds the current bounds, the node can be
pruned early.

Quiescence search significantly improves evaluation stability in tactical
positions and has become a standard component of classical chess engines. A
comprehensive discussion of the horizon effect and quiescence search can be
found in survey articles on game tree pruning~\cite{marsland1986review}.

\paragraph{Alpha-Beta Pruning}
Alpha-Beta pruning is an optimization technique that reduces the number of nodes
that need to be explored by Minimax or Negamax without affecting the correctness
of the final result. The technique was introduced by Knuth and Moore~\cite{knuth1975alphabeta}
and is a cornerstone of modern game tree search.

The algorithm maintains two bounds during search:
\begin{itemize}
    \item $\alpha$, the best score that the maximizing player can guarantee so
    far,
    \item $\beta$, the best score that the minimizing player can guarantee so
    far.
\end{itemize}

At the root of the search tree, $\alpha$ is initialized to $-\infty$ and $\beta$
to $+\infty$. As the search progresses, these bounds are updated based on the
values of explored child nodes. Whenever a node is encountered for which it can
be shown that its value cannot improve the current result, the corresponding
subtree is pruned and not searched further.

More concretely, during a maximizing step, the value of $\alpha$ is updated
whenever a better move is found. If $\alpha$ becomes greater than or equal to
$\beta$, further exploration of sibling nodes is unnecessary, as the minimizing
player would avoid reaching this position. An analogous argument applies during
minimizing steps.

\paragraph{Effectiveness and Influencing Factors}
In the worst case, Alpha-Beta pruning explores the same number of nodes as plain
Minimax. However, with optimal move ordering, the effective branching factor is
significantly reduced, allowing the search to reach roughly twice the depth of
unpruned Minimax for the same computational effort.

The effectiveness of Alpha-Beta pruning therefore depends strongly on the order
in which moves are explored. Techniques such as heuristic move ordering,
transposition tables, and iterative deepening are commonly employed to improve
this order and thereby increase pruning efficiency. These aspects are discussed
in more detail in later sections of this report.

For a comprehensive discussion of Alpha-Beta pruning and its variants in the
context of game-playing programs, see Marsland~\cite{marsland1986review}.

\subsection{Heuristic Evaluation}

Most positions encountered during search are non-terminal and therefore cannot be
assigned an exact game-theoretic value. Instead, they must be evaluated
heuristically. An evaluation function assigns a numerical score to a position,
indicating its desirability from the perspective of the side to move.

Classical evaluation functions combine material considerations with positional
features such as piece activity, king safety, and pawn structure. These features
are typically handcrafted and encode domain knowledge about the game of chess.
The design of the evaluation function has a significant impact on the playing
strength of the engine, particularly when search depth is limited.

\subsection{Time-Constrained Search}

In practical settings, chess engines are required to operate under strict time
constraints. Rather than searching to a fixed depth, the engine must dynamically
adapt its search behavior to the available time budget.

\paragraph{Iterative Deepening}
Iterative deepening is a widely used technique to address this requirement. The
search is performed repeatedly with increasing depth limits, starting from a
shallow search. This approach ensures that a valid move is always available, even
if the search is interrupted prematurely.

\paragraph{Interaction with Pruning}
In addition to providing robustness under time constraints, iterative deepening
also improves the effectiveness of Alpha-Beta pruning. Information gathered in
earlier iterations, such as good move orderings, can be reused in deeper searches,
leading to a significant reduction in the number of explored nodes.

The interaction between heuristic evaluation, Alpha-Beta pruning, and
time-constrained iterative deepening forms the foundation of modern chess
engines and is a central aspect of this work.

\subsection{Universal Chess Interface (UCI)}
\label{subsec:uci}

In practice, chess engines are controlled through standardized communication
protocols. The most widely used protocol in modern computer chess is the
\emph{Universal Chess Interface} (UCI).

UCI specifies a text-based command interface that allows external tools to set up
positions, initiate searches, and receive search results. In particular, search
constraints such as time limits, depth limits, node limits, and pondering are
communicated via the \texttt{go} command.

The protocol itself is independent of the engine’s internal implementation. As a
result, UCI naturally separates external control from internal search logic.
This separation is reflected in the design of the engine presented in this work,
where UCI parameters are first mapped to an abstract constraint representation
before being enforced during search.

All experiments reported in this work are conducted using UCI-compatible
interfaces and tournament frameworks. UCI therefore provides the practical
context in which time-constrained search and iterative deepening are evaluated.